% ============================================================
%  CHAPITRE III : Conception, réalisation et résultats
% ============================================================
\chapter{Conception, réalisation et résultats}
\label{chap:conception}

\section*{Introduction}

Ce chapitre décrit la mise en \oe uvre de la démarche. Il présente
successivement l'architecture de la solution, la chaîne de préparation des
données, l'analyse du portefeuille de références, le calcul des paramètres de
gestion, le tableau de bord développé, puis les résultats obtenus et leur
interprétation.

% -------------------------------------------------------
\section{Architecture de la solution}
\label{sec:architecture}

\subsection{Vue d'ensemble}

La figure~\ref{fig:architecture} illustre l'enchaînement des traitements, de
l'export brut fourni par l'\gls{ocp} jusqu'aux livrables destinés aux équipes.

\begin{figure}[H]
  \centering
  \begin{tikzpicture}[
    node distance=1.1cm and 1.6cm,
    block/.style={
      rectangle, draw=iacsTeal, fill=iacsLight!50,
      rounded corners=5pt, minimum width=2.9cm,
      minimum height=1.1cm, align=center,
      font=\scriptsize\sffamily, line width=0.8pt
    },
    dbblock/.style={
      cylinder, draw=iacsTeal, fill=iacsGreen!15,
      shape border rotate=90, aspect=0.25,
      minimum width=1.9cm, minimum height=1.1cm,
      font=\scriptsize\sffamily, line width=0.8pt
    },
    outblock/.style={block, fill=iacsOrange!15, draw=iacsOrange},
    arrow/.style={-Stealth, thick, color=iacsTeal},
    dataarrow/.style={-Stealth, thick, color=iacsAccent, dashed}
  ]
    % Couche données
    \node[dbblock] (source) {Export\\Excel};
    \node[block, right=of source] (clean) {Nettoyage\\\& contrôles};
    \node[block, right=of clean]  (feat)  {Calcul des\\indicateurs};

    % Couche analyse
    \node[block, below=1.4cm of feat, fill=iacsTeal!15]
      (segment) {\textbf{Segmentation}\\ABC / XYZ};
    \node[block, left=of segment, fill=iacsTeal!15]
      (calc) {\textbf{Paramètres}\\SS, ROP, EOQ};
    \node[block, left=of calc, fill=iacsTeal!15]
      (simul) {\textbf{Simulation}\\service / coût};

    % Couche restitution
    \node[outblock, below=1.4cm of calc] (dash) {Tableau de bord\\Streamlit};
    \node[outblock, left=of dash]  (xls)  {Export\\Excel};
    \node[outblock, right=of dash] (rap)  {Rapport\\\& graphiques};

    \draw[arrow] (source) -- (clean);
    \draw[arrow] (clean) -- (feat);
    \draw[arrow] (feat) -- (segment);
    \draw[arrow] (segment) -- (calc);
    \draw[arrow] (calc) -- (simul);
    \draw[arrow] (calc) -- (dash);
    \draw[dataarrow] (simul) -- (xls);
    \draw[dataarrow] (segment.south) to[out=-60, in=30] (rap.north east);
  \end{tikzpicture}
  \caption{Architecture générale de la solution}
  \label{fig:architecture}
\end{figure}

\subsection{Choix technologiques}

L'ensemble de la solution repose sur des outils libres et gratuits, exécutés
sous Debian dans un environnement virtuel Python isolé. Le
tableau~\ref{tab:outils} de l'annexe~A détaille le rôle de chaque composant.

\iacsnote{Le choix d'une pile entièrement libre répond à une contrainte du
cahier des charges : la solution doit rester reproductible et transférable sans
coût de licence.}

% -------------------------------------------------------
\section{Préparation des données}
\label{sec:preparation}

\subsection{Sources et périmètre}

\hspace{1cm}[Décrire l'export reçu : nature, période couverte, nombre de
références, nombre de lignes de mouvement, périmètre retenu.]

\subsection{Contrôles de qualité}

Les contrôles appliqués avant toute analyse sont résumés dans le
tableau~\ref{tab:controles}.

\begin{table}[H]
  \centering
  \caption{Contrôles de qualité appliqués aux données}
  \label{tab:controles}
  \renewcommand{\arraystretch}{1.4}
  \begin{tabularx}{\textwidth}{|l|X|c|}
    \hline
    \rowcolor{iacsTableHead}
    \textcolor{white}{\textbf{Contrôle}} &
    \textcolor{white}{\textbf{Traitement appliqué}} &
    \textcolor{white}{\textbf{Lignes concernées}} \\
    \hline
    Doublons            & Suppression des mouvements identiques        & [\ldots] \\
    \hline
    \rowcolor{iacsLightGray}
    Valeurs manquantes  & Imputation ou exclusion selon le champ       & [\ldots] \\
    \hline
    Quantités négatives & Distinction retour / erreur de saisie        & [\ldots] \\
    \hline
    \rowcolor{iacsLightGray}
    Dates aberrantes    & Exclusion hors période d'étude               & [\ldots] \\
    \hline
    Références inactives& Isolement des articles sans mouvement        & [\ldots] \\
    \hline
  \end{tabularx}
\end{table}

\subsection{Construction de la série de consommation}

L'algorithme~\ref{algo:agregation} décrit la construction de la série
mensuelle de consommation par référence.

\begin{algorithm}[H]
  \caption{Agrégation de la consommation par référence et par mois}
  \label{algo:agregation}
  \KwEntree{Table des mouvements de stock $M$, période d'étude $[t_{0}, t_{1}]$}
  \KwSortie{Matrice de consommation mensuelle $C$}
  Filtrer $M$ sur les mouvements de sortie et sur $[t_{0}, t_{1}]$\;
  Extraire le mois de chaque mouvement\;
  \ForEach{référence $i$}{
    \ForEach{mois $t$ de la période}{
      $C[i, t] \leftarrow$ somme des quantités sorties\;
      \If{aucun mouvement}{
        $C[i, t] \leftarrow 0$ \tcp*{mois sans consommation}
      }
    }
  }
  \KwRet{$C$}\;
\end{algorithm}

\begin{lstlisting}[language=Python, caption={Agrégation de la consommation mensuelle}, label={lst:agregation}]
import pandas as pd

mouvements = pd.read_excel("data/raw/mouvements.xlsx")

sorties = mouvements[mouvements["type_mouvement"] == "SORTIE"].copy()
sorties["mois"] = sorties["date_mouvement"].dt.to_period("M")

conso = (
    sorties
    .groupby(["code_article", "mois"], as_index=False)["quantite"]
    .sum()
    .pivot(index="code_article", columns="mois", values="quantite")
    .fillna(0)
)
\end{lstlisting}

% -------------------------------------------------------
\section{Analyse du portefeuille de références}
\label{sec:analyse_portefeuille}

\subsection{Classification ABC}

\hspace{1cm}[Insérer la courbe de Pareto obtenue et commenter la répartition
réelle entre les classes A, B et C.]

\begin{table}[H]
  \centering
  \caption{Résultats de la classification ABC}
  \label{tab:resultats_abc}
  \renewcommand{\arraystretch}{1.4}
  \begin{tabularx}{\textwidth}{|c|c|c|c|X|}
    \hline
    \rowcolor{iacsTableHead}
    \textcolor{white}{\textbf{Classe}} &
    \textcolor{white}{\textbf{Nb réf.}} &
    \textcolor{white}{\textbf{\% réf.}} &
    \textcolor{white}{\textbf{\% valeur}} &
    \textcolor{white}{\textbf{Commentaire}} \\
    \hline
    A & [\ldots] & [\ldots] & [\ldots] & \\
    \hline
    \rowcolor{iacsLightGray}
    B & [\ldots] & [\ldots] & [\ldots] & \\
    \hline
    C & [\ldots] & [\ldots] & [\ldots] & \\
    \hline
  \end{tabularx}
\end{table}

\subsection{Classification XYZ et matrice croisée}

\hspace{1cm}[Présenter la distribution du coefficient de variation, les seuils
retenus pour séparer X, Y et Z, puis la matrice croisée ABC/XYZ à neuf cases.]

\subsection{Analyse des délais fournisseurs}

\hspace{1cm}[Comparer les délais théoriques et les délais réellement observés,
et caractériser leur dispersion $\sigma_{L}$ par famille ou par fournisseur.]

% -------------------------------------------------------
\section{Calcul des paramètres de gestion}
\label{sec:calcul_parametres}

\subsection{Règles retenues par segment}

Le tableau~\ref{tab:regles_segment} associe à chaque case de la matrice
ABC/XYZ un mode de gestion et un taux de service cible.

\begin{table}[H]
  \centering
  \caption{Règles de gestion par segment}
  \label{tab:regles_segment}
  \renewcommand{\arraystretch}{1.4}
  \begin{tabularx}{\textwidth}{|c|X|c|}
    \hline
    \rowcolor{iacsTableHead}
    \textcolor{white}{\textbf{Segment}} &
    \textcolor{white}{\textbf{Mode de gestion}} &
    \textcolor{white}{\textbf{Service cible}} \\
    \hline
    AX & Point de commande automatique, suivi hebdomadaire & [\ldots] \\
    \hline
    \rowcolor{iacsLightGray}
    AY / AZ & Suivi rapproché, arbitrage manuel & [\ldots] \\
    \hline
    BX / BY & Règles automatiques, revue mensuelle & [\ldots] \\
    \hline
    \rowcolor{iacsLightGray}
    CX / CY & Recomplètement périodique, lots larges & [\ldots] \\
    \hline
    BZ / CZ & Gestion à la demande, pas de stock permanent & [\ldots] \\
    \hline
  \end{tabularx}
\end{table}

\subsection{Implémentation}

\begin{lstlisting}[language=Python, caption={Calcul du stock de sécurité et du point de commande}, label={lst:ss_rop}]
import numpy as np
from scipy.stats import norm

def stock_securite(sigma_d, d_moyen, lead_time, sigma_lt, service_cible):
    """Stock de securite avec demande et delai incertains."""
    z = norm.ppf(service_cible)
    return z * np.sqrt(lead_time * sigma_d**2 + d_moyen**2 * sigma_lt**2)

def point_de_commande(d_moyen, lead_time, ss):
    return d_moyen * lead_time + ss

def quantite_economique(demande_annuelle, cout_passation, cout_possession):
    return np.sqrt(2 * demande_annuelle * cout_passation / cout_possession)
\end{lstlisting}

\subsection{Hypothèses de coûts}

\iacswarning{Le coût de passation $S$ et le taux de possession $i$ ne figurent
pas dans l'export de données ; ils doivent être estimés avec les services
concernés. Les résultats seront présentés sous forme d'analyse de sensibilité à
ces deux paramètres.}

% -------------------------------------------------------
\section{Tableau de bord}
\label{sec:dashboard}

\subsection{Fonctionnalités}

\begin{itemize}
  \item Vue d'ensemble du portefeuille : valeur du stock, taux de rotation,
    part des références dormantes.
  \item Exploration par famille, par classe \gls{abc}/\gls{xyz} ou par
    fournisseur.
  \item Fiche par référence : historique de consommation, paramètres calculés,
    couverture actuelle.
  \item Simulateur : effet d'un taux de service cible sur la valeur totale du
    stock de sécurité.
  \item Export des résultats au format Excel.
\end{itemize}

\subsection{Captures d'écran}

\hspace{1cm}[Insérer ici les captures du tableau de bord avec
\texttt{\textbackslash includegraphics}, en plaçant les images dans
\texttt{figures/}.]

% -------------------------------------------------------
\section{Résultats}
\label{sec:resultats}

\subsection{Situation actuelle et situation optimisée}

\begin{table}[H]
  \centering
  \caption{Comparaison avant / après optimisation}
  \label{tab:comparaison}
  \renewcommand{\arraystretch}{1.4}
  \begin{tabularx}{\textwidth}{|X|c|c|c|}
    \hline
    \rowcolor{iacsTableHead}
    \textcolor{white}{\textbf{Indicateur}} &
    \textcolor{white}{\textbf{Actuel}} &
    \textcolor{white}{\textbf{Optimisé}} &
    \textcolor{white}{\textbf{Écart}} \\
    \hline
    Valeur du stock          & [\ldots] & [\ldots] & [\ldots] \\
    \hline
    \rowcolor{iacsLightGray}
    Taux de service estimé   & [\ldots] & [\ldots] & [\ldots] \\
    \hline
    Taux de rotation         & [\ldots] & [\ldots] & [\ldots] \\
    \hline
    \rowcolor{iacsLightGray}
    Couverture moyenne (j)   & [\ldots] & [\ldots] & [\ldots] \\
    \hline
    Valeur du stock dormant  & [\ldots] & [\ldots] & [\ldots] \\
    \hline
  \end{tabularx}
\end{table}

\subsection{Courbe d'arbitrage service / coût}

\hspace{1cm}[Insérer la courbe donnant la valeur du stock de sécurité en
fonction du taux de service cible, et identifier le point d'inflexion au-delà
duquel chaque point de service supplémentaire devient disproportionnellement
coûteux.]

\subsection{Analyse de sensibilité}

\hspace{1cm}[Présenter l'effet d'une variation du taux de possession et du coût
de passation sur les quantités de commande et sur les gains annoncés.]

% -------------------------------------------------------
\section{Discussion}
\label{sec:discussion}

\subsection{Interprétation}

\hspace{1cm}[Interpréter les résultats : quels segments concentrent les gains,
quelle est la part attribuable à la réduction du surstock et celle attribuable
à la fiabilisation des délais.]

\subsection{Limites}

\begin{itemize}
  \item Qualité et profondeur de l'historique disponible.
  \item Hypothèses de normalité peu adaptées aux demandes intermittentes.
  \item Coûts de passation, de possession et de rupture estimés et non mesurés.
  \item Résultats issus d'une simulation et non d'une mise en production.
\end{itemize}

\section*{Conclusion}

Ce chapitre a présenté la chaîne de traitement mise en place, les paramètres de
gestion calculés et les gains simulés. La conclusion générale dresse le bilan du
travail et ouvre sur les perspectives d'industrialisation.
